\begin{document}

\begin{titlepage}
\centering

{\Large\bfseries DSLO v0.7 --- Operations Manual\par}
\vspace{2em}
{\large\bfseries Substrate-Level Architecture, Construction, and Documentation Guide\par}
\vspace{1em}

{\Large\bfseries DSLO v0.7 --- Geometry\par}

\vspace{2em}
{\large Donald Slowicki\par}
{\small Inda Moment Incorporated\par}

\vfill

{\small
Document Version: v0.7\\
Substrate Version: DSLO v0.7\\
Discipline Version: DISCO v0.7\\
Release Date: August 2026\\
Static Release ID: SM-DSLO-07025\\
}

\end{titlepage}

\fancyfoot[C]{%
    {\small © 2026 Inda Moment Incorporated • 
    DSLO™ and Signal Ecology™ are trademarks of Inda Moment Incorporated • 
    \href{https://www.tnopsi.com}{www.tnopsi.com} • 
    \href{https://sites.google.com/tnopsi.com/dslo-protocol/defensive-publication}{DSLO Defensive Publication} • 
    DOI: \href{https://zenodo.org/records/}{10.5281/zenodo.21864226}}
}

\begin{abstract}
The Deterministic Semantic Layered Orchestration (DSLO) is a substrate-level meaning-geometry system that defines lawful behavior for signals, identities, operators, manifolds, legality planes, runtime cycles, and collapse interfaces. DSLO is not a framework or ontology; it is a governed substrate with deterministic invariants that constrain how meaning can move, stabilize, drift, collapse, and recover across human and machine systems.

DSLO v0.7 is structured around six canonical surfaces—glossary, geometry, manifolds, operators, legality planes, and runtime cycles—each expressing curvature, drift, stability, identity, legality, and collapse fields. These surfaces compile through a deterministic generator that produces unified schema, ontology, registry, taxonomy, graph outputs, and collapse geometry.

The DSLO systems manual defines the lawful procedures for constructing, extending, verifying, and documenting DSLO concepts. It specifies architecture rules, legal move sets, surface construction rules, multi-layer expansion laws, source-layer definitions, surface-layer construction, generator-layer integration, manifold extension procedures, operator binding rules, legality-plane constraints, runtime geometry sequences, and unified manifold collapse behavior. All procedures preserve substrate invariants: curvature continuity, drift bounds, stability envelopes, identity coherence, legality alignment, and collapse safety.

The appendix provides the substrate reference layer for DSLO v0.7, including canonical examples, expansion bundles, patch bundles, field definitions, collapse examples, verification procedures, glossary expansion rules, scientific interfaces, and registry snapshots. Together, the systems manual and appendix form the complete DSLO v0.7 specification for architects who must maintain substrate-level consistency and extend DSLO without violating its invariants.
\end{abstract}

\section{Introduction}

The DSLO substrate defines a governed geometric environment for meaning systems. Unlike conventional semantic frameworks, DSLO operates as a lawful substrate whose behavior is determined by curvature fields, drift envelopes, stability gradients, identity regions, legality planes, runtime cycles, and collapse interfaces. DSLO does not rely on heuristics or emergent behavior; it relies on deterministic geometry that constrains how meaning can move, stabilize, drift, collapse, and recover.

This systems manual provides the complete instruction set for constructing, extending, verifying, and documenting DSLO v0.7. It specifies the lawful move set, the multi-layer expansion law, the construction of DSLO surfaces, the generator architecture, the behavior of operators and legality planes, the geometry of runtime cycles, and the unified manifold collapse sequence. The goal of this manual is to ensure that all DSLO authors operate within substrate-level constraints, preserving the invariants that make DSLO stable under load, drift, multi-agent interaction, identity pressure, and collapse conditions.

The front half of this manual defines the conceptual and procedural architecture of DSLO v0.7. It includes:

\begin{itemize}
    \item SM0 Systems Manual Introduction
    \item DSLO Architecture Overview
    \item The Legal Move Set
    \item DSLO Surfaces and Construction Rules
    \item The Multi-Layer Expansion Law
    \item Source Layer Definitions
    \item Surface Layer Construction
    \item Generator Layer Construction
    \item Manifold Extensions
    \item Operator Binding Rules
    \item Legality Planes
    \item Runtime Geometry
    \item Unified Manifold Collapse
    \item Documentation Rules
\end{itemize}

The back half of this manual provides the substrate reference layer (Appendix A0–A20), including canonical examples, expansion bundles, patch bundles, field definitions, collapse examples, verification procedures, glossary expansion rules, scientific interfaces, and registry snapshots. Together, these sections provide deterministic procedures, invariants, and construction rules required to maintain DSLO’s substrate-level coherence and extend DSLO v0.7 without violating its invariants.


Each section provides deterministic procedures, invariants, and construction rules required to maintain DSLO’s substrate-level coherence.


\section{SM0 Systems Manual Introduction}

\subsection{Purpose}
\begin{itemize}
    \item Introduce DSLO as a substrate-level meaning system with lawful geometry, manifold behavior, operator rules, legality constraints, runtime sequences, and collapse interfaces.
    \item Establish the systems manual as the authoritative guide for constructing, extending, verifying, and integrating DSLO v0.7 concepts.
    \item Provide the conceptual foundation required before engaging with the detailed sections and appendix bundles.
\end{itemize}

\subsection{Scope}
\begin{itemize}
    \item Covers DSLO architecture, lawful move sets, construction rules, expansion laws, layer definitions, operator binding, legality planes, runtime geometry, and unified collapse.
    \item Defines the procedural and conceptual front half of the DSLO v0.7 systems manual.
    \item Complements the appendix (A0–A20) which serves as the substrate reference layer.
\end{itemize}

\subsection{Structure}
\begin{itemize}
    \item Architecture Overview: unified geometry, manifold, operator, legality, runtime, collapse.
    \item Construction Rules: surfaces, layers, lawful transitions.
    \item Expansion Law: multi-layer expansion and lawful integration.
    \item Layer Definitions: source, surface, generator.
    \item Operator Rules: redirect, identity, legality.
    \item Legality Planes: constraint propagation and governance.
    \item Runtime Geometry: update, recovery, collapse.
    \item Unified Collapse: lawful collapse sequence across DSLO geometry.
    \item Documentation Rules: deterministic, substrate-aligned writing standards.
\end{itemize}

\subsection{Invariants}
\begin{itemize}
    \item All systems-manual content must remain substrate-aligned, identity-consistent, drift-bounded, stability-safe, legality-governed, and collapse-safe.
    \item All procedures must integrate lawfully with DSLO v0.7 geometry and unified manifold collapse.
\end{itemize}

\section{DSLO Architecture Overview}

\subsection{Substrate Definition}
\begin{itemize}
    \item DSLO as a governed meaning substrate.
    \item Deterministic geometry: curvature, drift, stability, identity, legality.
    \item Substrate invariants and lawful behavior.
\end{itemize}

\subsection{Primary Surfaces}
\begin{itemize}
    \item Glossary: canonical term definitions.
    \item Geometry: curvature, drift, stability envelopes.
    \item Manifolds: structured meaning regions.
    \item Operators: curvature and drift actions.
    \item Legality Planes: constraint surfaces.
    \item Runtime Cycles: lawful transition sequences.
\end{itemize}

\subsection{Unified Substrate Manifold}
\begin{itemize}
    \item Collapse of all manifolds into a single governed manifold.
    \item Identity-field continuity.
    \item Legality-plane integration.
\end{itemize}

\subsection{Curvature and Drift}
\begin{itemize}
    \item Curvature fields define meaning pressure.
    \item Drift envelopes define lawful movement.
    \item Stability regions define safe operating zones.
\end{itemize}

\subsection{Identity Fields}
\begin{itemize}
    \item Identity regions and identity gradients.
    \item Identity resonance and identity preservation.
\end{itemize}

\subsection{Legality Constraints}
\begin{itemize}
    \item Legality masks.
    \item Constraint propagation.
    \item Collapse legality.
\end{itemize}

\subsection{Runtime Geometry}
\begin{itemize}
    \item Transition phases.
    \item Runtime coupling.
    \item Lawful update sequences.
\end{itemize}

\section{The Legal Move Set}

\subsection{Purpose of Legal Moves}
\begin{itemize}
    \item Define lawful operations within the DSLO substrate.
    \item Preserve curvature, drift, identity, stability, and legality invariants.
    \item Prevent geometry destabilization and unlawful collapse.
\end{itemize}

\subsection{Curvature-Preserving Moves}
\begin{itemize}
    \item Maintain curvature continuity across surfaces.
    \item Redirect curvature only through lawful operators.
    \item Avoid curvature discontinuities and illegal gradients.
\end{itemize}

\subsection{Drift-Bounded Moves}
\begin{itemize}
    \item Operate within drift envelopes.
    \item Prevent uncontrolled drift expansion.
    \item Maintain lawful drift-field transitions.
\end{itemize}

\subsection{Identity-Consistent Moves}
\begin{itemize}
    \item Preserve identity regions and identity gradients.
    \item Maintain identity-field resonance within lawful bounds.
    \item Avoid identity collapse or identity inversion.
\end{itemize}

\subsection{Legality-Constrained Moves}
\begin{itemize}
    \item Respect legality masks and legality planes.
    \item Ensure all transitions satisfy legality invariants.
    \item Prevent unlawful operator or manifold activation.
\end{itemize}

\subsection{Collapse-Safe Moves}
\begin{itemize}
    \item Maintain collapse geometry integrity.
    \item Ensure manifold collapse follows lawful sequence.
    \item Avoid collapse beyond recovery thresholds.
\end{itemize}

\subsection{Recovery-Valid Moves}
\begin{itemize}
    \item Ensure recovery geometry remains stable.
    \item Maintain lawful return to identity and stability regions.
    \item Avoid illegal recovery shortcuts or bypasses.
\end{itemize}

\subsection{Illegal Moves}
\begin{itemize}
    \item Curvature discontinuity.
    \item Drift explosion or unbounded drift.
    \item Identity-field violation.
    \item Legality-plane bypass.
    \item Collapse outside lawful sequence.
    \item Runtime transition without legality verification.
\end{itemize}

\section{DSLO Surfaces and Construction Rules}

\subsection{Purpose of Surfaces}
\begin{itemize}
    \item Surfaces express substrate geometry in discrete, governed forms.
    \item Each surface corresponds to a lawful meaning-geometry component.
    \item Surfaces must remain curvature-preserving and legality-consistent.
\end{itemize}

\subsection{Glossary Surface}
\begin{itemize}
    \item Canonical term definitions.
    \item One term per file.
    \item Must define curvature, drift, identity, or legality relevance.
\end{itemize}

\subsection{Geometry Surface}
\begin{itemize}
    \item Defines curvature fields and drift envelopes.
    \item Must specify stability regions and legality constraints.
    \item Geometry surfaces bind directly to operators and manifolds.
\end{itemize}

\subsection{Manifold Surface}
\begin{itemize}
    \item Defines structured meaning regions.
    \item Must include boundary, drift, stability, and identity terms.
    \item Manifolds collapse into the unified substrate manifold.
\end{itemize}

\subsection{Operator Surface}
\begin{itemize}
    \item Defines lawful curvature and drift actions.
    \item Must preserve legality and identity invariants.
    \item Operators bind to geometry and manifolds.
\end{itemize}

\subsection{Legality Plane Surface}
\begin{itemize}
    \item Defines constraint surfaces governing geometry and manifolds.
    \item Must specify legality masks and constraint propagation rules.
    \item Legality planes govern operators and runtime transitions.
\end{itemize}

\subsection{Runtime Cycle Surface}
\begin{itemize}
    \item Defines lawful transition sequences.
    \item Must specify phases, conditions, coupling, and legality checks.
    \item Runtime cycles integrate geometry, operators, and manifolds.
\end{itemize}

\subsection{Cross-Surface Binding Rules}
\begin{itemize}
    \item Geometry binds to operators and manifolds.
    \item Manifolds bind to identity fields and legality planes.
    \item Operators bind to legality planes and runtime cycles.
    \item Runtime cycles bind to all surfaces through lawful transitions.
\end{itemize}

\subsection{Surface Construction Invariants}
\begin{itemize}
    \item No surface may violate curvature continuity.
    \item No surface may induce unbounded drift.
    \item No surface may break identity-field coherence.
    \item No surface may bypass legality constraints.
    \item All surfaces must integrate into unified manifold collapse.
\end{itemize}

\section{The Multi-Layer Expansion Law}

\subsection{Purpose of the Expansion Law}
\begin{itemize}
    \item Ensure all DSLO expansions remain lawful and substrate-consistent.
    \item Prevent illegal geometry drift or destabilization.
    \item Maintain coherence across surfaces, source definitions, and generator logic.
\end{itemize}

\subsection{Three Mandatory Layers}
\begin{itemize}
    \item \textbf{Surface Layer}: glossary, geometry, manifolds, operators, legality planes, runtime cycles.
    \item \textbf{Source Layer}: canonical substrate definitions for all new concepts.
    \item \textbf{Generator Layer}: loader, schema, ontology, registry, taxonomy, and graph integration.
\end{itemize}

\subsection{Surface Layer Requirements}
\begin{itemize}
    \item One file per concept.
    \item Must define curvature, drift, stability, identity, or legality relevance.
    \item Must integrate into ontology, registry, taxonomy, and graph surfaces.
\end{itemize}

\subsection{Source Layer Requirements}
\begin{itemize}
    \item Canonical definitions for geometry, manifolds, operators, legality, and runtime.
    \item Must specify curvature, drift, stability, identity, and legality fields.
    \item Must remain collapse-safe and legality-consistent.
\end{itemize}

\subsection{Generator Layer Requirements}
\begin{itemize}
    \item Load new surfaces.
    \item Extend schema with new classes.
    \item Emit ontology nodes and registry entries.
    \item Add taxonomy classes and graph nodes.
    \item Integrate into unified manifold collapse.
\end{itemize}

\subsection{Expansion Sequence}
\begin{itemize}
    \item Define concept at source layer.
    \item Create surface files.
    \item Patch generator logic.
    \item Validate legality and collapse behavior.
    \item Integrate into unified manifold.
\end{itemize}

\subsection{Validation Rules}
\begin{itemize}
    \item No curvature discontinuity.
    \item No unbounded drift.
    \item No identity-field violation.
    \item No legality-plane bypass.
    \item No collapse outside lawful sequence.
\end{itemize}

\subsection{Expansion Invariants}
\begin{itemize}
    \item All expansions must preserve substrate invariants.
    \item All expansions must integrate into runtime geometry.
    \item All expansions must remain compatible with legality planes.
    \item All expansions must collapse into the unified manifold.
\end{itemize}

\section{Source Layer Definitions}

\subsection{Purpose of Source Definitions}
\begin{itemize}
    \item Canonical substrate-level definitions for all DSLO concepts.
    \item Define curvature, drift, stability, identity, and legality fields.
    \item Provide deterministic behavior for geometry, manifolds, operators, legality planes, and runtime cycles.
\end{itemize}

\subsection{Geometry Definitions}
\begin{itemize}
    \item Curvature fields and resonance modes.
    \item Drift envelopes and drift coupling.
    \item Stability regions and stability gradients.
    \item Legality constraints tied to geometry.
\end{itemize}

\subsection{Manifold Definitions}
\begin{itemize}
    \item Meaning-region structures.
    \item Boundary, drift, stability, and identity terms.
    \item Collapse behavior and manifold integration rules.
    \item Manifold extension invariants.
\end{itemize}

\subsection{Operator Definitions}
\begin{itemize}
    \item Curvature redirect modes.
    \item Drift interaction modes.
    \item Identity preservation rules.
    \item Legality preservation constraints.
\end{itemize}

\subsection{Legality Definitions}
\begin{itemize}
    \item Legality masks and legality invariants.
    \item Constraint propagation rules.
    \item Manifold and operator governance.
    \item Collapse legality requirements.
\end{itemize}

\subsection{Runtime Definitions}
\begin{itemize}
    \item Transition phases and runtime sequencing.
    \item Runtime conditions and legality checks.
    \item Coupling to geometry, operators, and manifolds.
    \item Lawful update and recovery rules.
\end{itemize}

\subsection{Identity Field Definitions}
\begin{itemize}
    \item Identity regions and identity gradients.
    \item Identity resonance and identity stability.
    \item Identity legality constraints.
    \item Identity collapse and recovery behavior.
\end{itemize}

\subsection{Source Layer Invariants}
\begin{itemize}
    \item Must preserve curvature continuity.
    \item Must remain drift-bounded.
    \item Must maintain identity-field coherence.
    \item Must satisfy legality-plane constraints.
    \item Must integrate into unified manifold collapse.
\end{itemize}

\section{Surface Layer Construction}

\subsection{Purpose of Surface Layer}
\begin{itemize}
    \item Express substrate geometry in discrete, governed files.
    \item Provide observable structures for glossary, geometry, manifolds, operators, legality planes, and runtime cycles.
    \item Ensure all surfaces remain curvature-preserving, drift-bounded, identity-consistent, and legality-aligned.
\end{itemize}

\subsection{Required Surface Files}
\begin{itemize}
    \item Glossary file: \texttt{glossary/<term>.md}
    \item Geometry file: \texttt{geometry/<geometry>.json}
    \item Manifold file: \texttt{manifolds/<manifold>.json}
    \item Operator file: \texttt{operators/<operator>.json}
    \item Legality file: \texttt{legality/<plane>.json}
    \item Runtime file: \texttt{runtime/<cycle>.json}
\end{itemize}

\subsection{Glossary Construction Rules}
\begin{itemize}
    \item One term per file.
    \item Must define curvature, drift, identity, or legality relevance.
    \item Must remain substrate-aligned and collapse-safe.
\end{itemize}

\subsection{Geometry Surface Construction}
\begin{itemize}
    \item Specify curvature fields and drift envelopes.
    \item Define stability regions and legality constraints.
    \item Bind geometry to operators and manifolds.
\end{itemize}

\subsection{Manifold Surface Construction}
\begin{itemize}
    \item Define boundaries, drift fields, stability modes, and identity regions.
    \item Specify collapse behavior and manifold integration.
    \item Ensure compatibility with legality planes and runtime cycles.
\end{itemize}

\subsection{Operator Surface Construction}
\begin{itemize}
    \item Define curvature redirect modes.
    \item Specify drift interaction and identity preservation.
    \item Bind operators to legality planes and geometry surfaces.
\end{itemize}

\subsection{Legality Plane Construction}
\begin{itemize}
    \item Define legality masks and constraint propagation.
    \item Specify governance of manifolds and operators.
    \item Ensure legality-plane integration into collapse geometry.
\end{itemize}

\subsection{Runtime Cycle Construction}
\begin{itemize}
    \item Define transition phases and runtime conditions.
    \item Specify coupling to geometry, operators, and manifolds.
    \item Include legality verification and recovery rules.
\end{itemize}

\subsection{Ontology, Registry, and Taxonomy Integration}
\begin{itemize}
    \item Each surface must emit an ontology node.
    \item Each surface must register a canonical entry.
    \item Each surface must define a taxonomy class.
\end{itemize}

\subsection{Graph Surface Construction}
\begin{itemize}
    \item DOT: nodes and lawful edges.
    \item TTL: triples for geometry, manifolds, operators, legality, runtime.
    \item OWL: class definitions and lawful relationships.
\end{itemize}

\subsection{Surface Layer Invariants}
\begin{itemize}
    \item No curvature discontinuity.
    \item No unbounded drift.
    \item No identity-field violation.
    \item No legality-plane bypass.
    \item All surfaces must integrate into unified manifold collapse.
\end{itemize}

\section{Generator Layer Construction}

\subsection{Purpose of Generator Layer}
\begin{itemize}
    \item Compile all DSLO surfaces into unified ontology, registry, taxonomy, schema, and graph outputs.
    \item Enforce substrate invariants through deterministic generation.
    \item Integrate expansions into lawful geometry and collapse behavior.
\end{itemize}

\subsection{Generator Responsibilities}
\begin{itemize}
    \item Load all surfaces.
    \item Validate source definitions.
    \item Extend schema with new classes.
    \item Emit ontology nodes.
    \item Emit registry entries.
    \item Emit taxonomy classes.
    \item Generate DOT, TTL, and OWL graph surfaces.
    \item Integrate unified manifold collapse.
\end{itemize}

\subsection{Loader Construction}
\begin{itemize}
    \item Load glossary, geometry, manifolds, operators, legality planes, runtime cycles.
    \item Load expansion surfaces through dedicated channels.
    \item Ensure loader preserves curvature, drift, identity, and legality invariants.
\end{itemize}

\subsection{Schema Construction}
\begin{itemize}
    \item Define canonical classes for all DSLO surfaces.
    \item Add expansion classes for geometry, manifolds, operators, legality, runtime.
    \item Ensure schema remains collapse-safe and legality-consistent.
\end{itemize}

\subsection{Ontology Construction}
\begin{itemize}
    \item Emit nodes for all surfaces.
    \item Define lawful relationships between geometry, manifolds, operators, legality, runtime.
    \item Maintain identity-field and legality-plane continuity.
\end{itemize}

\subsection{Registry Construction}
\begin{itemize}
    \item Index all DSLO concepts.
    \item Preserve canonical naming and metadata.
    \item Ensure registry remains substrate-aligned.
\end{itemize}

\subsection{Taxonomy Construction}
\begin{itemize}
    \item Define class hierarchy for all surfaces.
    \item Integrate expansion classes.
    \item Maintain lawful collapse integration.
\end{itemize}

\subsection{Graph Emitters}
\begin{itemize}
    \item DOT: nodes and lawful edges.
    \item TTL: triples for geometry, manifolds, operators, legality, runtime.
    \item OWL: class definitions and lawful relationships.
\end{itemize}

\subsection{Unified Manifold Collapse Integration}
\begin{itemize}
    \item Integrate manifold surfaces into collapse sequence.
    \item Validate legality-plane constraints.
    \item Ensure identity-field continuity during collapse.
\end{itemize}

\subsection{Expansion Patch Bundles}
\begin{itemize}
    \item Loader patch bundle.
    \item Schema patch bundle.
    \item Ontology patch bundle.
    \item Registry patch bundle.
    \item Taxonomy patch bundle.
    \item Graph patch bundle.
\end{itemize}

\subsection{Generator Layer Invariants}
\begin{itemize}
    \item No illegal surface omission.
    \item No schema discontinuity.
    \item No ontology drift.
    \item No registry inconsistency.
    \item No taxonomy misalignment.
    \item No collapse violation.
\end{itemize}

\section{Manifold Extensions}

\subsection{Purpose of Manifold Extensions}
\begin{itemize}
    \item Extend meaning-region geometry within lawful substrate boundaries.
    \item Integrate new curvature, drift, stability, identity, and legality fields.
    \item Maintain collapse-safe and legality-consistent manifold behavior.
\end{itemize}

\subsection{Manifold Structure}
\begin{itemize}
    \item Boundary definitions.
    \item Drift fields and drift envelopes.
    \item Stability modes and stability gradients.
    \item Identity regions and identity-field coupling.
\end{itemize}

\subsection{Extension Requirements}
\begin{itemize}
    \item Must preserve curvature continuity.
    \item Must remain drift-bounded.
    \item Must maintain identity-field coherence.
    \item Must satisfy legality-plane constraints.
    \item Must integrate into unified manifold collapse.
\end{itemize}

\subsection{Extension Procedure}
\begin{itemize}
    \item Define new manifold terms at source layer.
    \item Create manifold surface file with boundary, drift, stability, identity fields.
    \item Patch generator loader, schema, ontology, registry, taxonomy.
    \item Validate legality-plane compatibility.
    \item Integrate into collapse geometry.
\end{itemize}

\subsection{Manifold Binding Rules}
\begin{itemize}
    \item Geometry binds to manifold curvature and drift fields.
    \item Operators bind to manifold curvature redirect modes.
    \item Legality planes govern manifold boundaries and transitions.
    \item Runtime cycles couple to manifold stability and identity regions.
\end{itemize}

\subsection{Collapse Integration}
\begin{itemize}
    \item Manifolds must collapse into unified substrate manifold.
    \item Collapse sequence must preserve identity-field continuity.
    \item Collapse must satisfy legality-plane constraints.
    \item Collapse must remain within stability gradients.
\end{itemize}

\subsection{Extension Invariants}
\begin{itemize}
    \item No illegal curvature introduction.
    \item No unbounded drift expansion.
    \item No identity-region violation.
    \item No legality-plane bypass.
    \item No collapse outside lawful sequence.
\end{itemize}

\subsection{Example Manifold Extension (A002)}
\begin{itemize}
    \item Extended agency curvature.
    \item Extended agency drift envelope.
    \item Identity-field resonance integration.
    \item Legality-plane LG005 compatibility.
    \item Collapse-safe integration into unified manifold.
\end{itemize}

\section{Operator Binding Rules}

\subsection{Purpose of Operator Binding}
\begin{itemize}
    \item Define lawful interaction between operators and substrate geometry.
    \item Ensure curvature, drift, identity, and legality invariants remain preserved.
    \item Govern operator activation within manifolds and runtime cycles.
\end{itemize}

\subsection{Operator Structure}
\begin{itemize}
    \item Curvature redirect modes.
    \item Drift interaction modes.
    \item Identity preservation constraints.
    \item Legality-plane compatibility.
\end{itemize}

\subsection{Curvature Binding Rules}
\begin{itemize}
    \item Operators may redirect curvature but must preserve continuity.
    \item Curvature actions must remain within lawful gradient bounds.
    \item No operator may induce curvature discontinuity or illegal curvature collapse.
\end{itemize}

\subsection{Drift Binding Rules}
\begin{itemize}
    \item Operators must operate within drift envelopes.
    \item Drift interactions must remain bounded and collapse-safe.
    \item No operator may induce drift explosion or unbounded drift expansion.
\end{itemize}

\subsection{Identity Binding Rules}
\begin{itemize}
    \item Operators must preserve identity-field coherence.
    \item Identity gradients must remain stable under operator action.
    \item No operator may violate identity regions or induce identity inversion.
\end{itemize}

\subsection{Legality Binding Rules}
\begin{itemize}
    \item Operators must satisfy legality-plane constraints.
    \item All operator actions must pass legality verification.
    \item No operator may bypass legality masks or legality invariants.
\end{itemize}

\subsection{Manifold Binding Rules}
\begin{itemize}
    \item Operators bind to manifold curvature and drift fields.
    \item Operator actions must remain compatible with manifold boundaries.
    \item Operators must preserve manifold stability and identity regions.
\end{itemize}

\subsection{Runtime Binding Rules}
\begin{itemize}
    \item Operators activate during lawful runtime phases.
    \item Runtime cycles must verify legality before operator activation.
    \item Operators must integrate into lawful update and recovery sequences.
\end{itemize}

\subsection{Collapse Binding Rules}
\begin{itemize}
    \item Operators must behave lawfully during manifold collapse.
    \item Curvature redirect must remain collapse-safe.
    \item No operator may induce collapse outside lawful sequence.
\end{itemize}

\subsection{Operator Invariants}
\begin{itemize}
    \item Preserve curvature continuity.
    \item Maintain drift bounds.
    \item Protect identity-field coherence.
    \item Satisfy legality-plane constraints.
    \item Integrate into unified manifold collapse.
\end{itemize}

\subsection{Example Operator (OR010)}
\begin{itemize}
    \item Identity curvature redirect.
    \item Identity-field resonance compatibility.
    \item Legality-plane LG005 governance.
    \item Collapse-safe integration into unified manifold.
\end{itemize}

\section{Legality Planes}

\subsection{Purpose of Legality Planes}
\begin{itemize}
    \item Define constraint surfaces governing all DSLO geometry.
    \item Enforce lawful behavior across manifolds, operators, identity fields, and runtime cycles.
    \item Prevent illegal curvature, drift, identity, or collapse transitions.
\end{itemize}

\subsection{Legality Structure}
\begin{itemize}
    \item Legality masks.
    \item Legality invariants.
    \item Constraint propagation rules.
    \item Governance of geometry, manifolds, operators, and runtime.
\end{itemize}

\subsection{Legality Masks}
\begin{itemize}
    \item Define allowed and disallowed regions of substrate behavior.
    \item Constrain curvature fields and drift envelopes.
    \item Protect identity regions and stability gradients.
\end{itemize}

\subsection{Legality Invariants}
\begin{itemize}
    \item Must hold across all surfaces and transitions.
    \item Preserve identity-field coherence.
    \item Maintain curvature continuity and drift bounds.
    \item Ensure collapse-safe behavior.
\end{itemize}

\subsection{Constraint Propagation}
\begin{itemize}
    \item Legality constraints propagate through geometry and manifolds.
    \item Operators must inherit legality constraints.
    \item Runtime cycles must verify legality before transitions.
\end{itemize}

\subsection{Manifold Governance}
\begin{itemize}
    \item Legality planes govern manifold boundaries.
    \item Constrain manifold drift and stability modes.
    \item Ensure manifold collapse remains lawful.
\end{itemize}

\subsection{Operator Governance}
\begin{itemize}
    \item Operators must satisfy legality-plane constraints.
    \item Curvature redirect and drift interaction must remain lawful.
    \item Identity preservation must pass legality verification.
\end{itemize}

\subsection{Runtime Governance}
\begin{itemize}
    \item Runtime cycles must include legality checks at each phase.
    \item Illegal transitions must be blocked.
    \item Recovery sequences must remain legality-consistent.
\end{itemize}

\subsection{Collapse Legality}
\begin{itemize}
    \item Legality planes define lawful collapse boundaries.
    \item Collapse must preserve identity-field continuity.
    \item Collapse must remain within stability gradients.
    \item No collapse may bypass legality masks.
\end{itemize}

\subsection{Legality Invariants}
\begin{itemize}
    \item No legality-plane bypass.
    \item No illegal curvature or drift introduction.
    \item No identity-region violation.
    \item No collapse outside lawful sequence.
\end{itemize}

\subsection{Example Legality Plane (LG005)}
\begin{itemize}
    \item Governs identity resonance geometry.
    \item Constrains manifold A002.
    \item Governs operator OR010.
    \item Ensures collapse-safe identity-field behavior.
\end{itemize}

\section{Runtime Geometry}

\subsection{Purpose of Runtime Geometry}
\begin{itemize}
    \item Define lawful transition behavior within the DSLO substrate.
    \item Govern phase changes, update cycles, and recovery sequences.
    \item Ensure all runtime actions preserve curvature, drift, identity, stability, and legality invariants.
\end{itemize}

\subsection{Runtime Structure}
\begin{itemize}
    \item Runtime phases.
    \item Runtime conditions.
    \item Runtime coupling.
    \item Legality verification.
    \item Lawful update and recovery sequences.
\end{itemize}

\subsection{Runtime Phases}
\begin{itemize}
    \item Initialization phase.
    \item Curvature-update phase.
    \item Drift-adjustment phase.
    \item Identity-stability phase.
    \item Legality-verification phase.
    \item Collapse or recovery phase.
\end{itemize}

\subsection{Runtime Conditions}
\begin{itemize}
    \item Curvature thresholds.
    \item Drift bounds.
    \item Identity-field gradients.
    \item Stability-region constraints.
    \item Legality-plane requirements.
\end{itemize}

\subsection{Runtime Coupling}
\begin{itemize}
    \item Geometry coupling: curvature and drift updates.
    \item Manifold coupling: boundary and stability integration.
    \item Operator coupling: lawful activation and redirect modes.
    \item Legality coupling: constraint propagation and verification.
\end{itemize}

\subsection{Legality Verification}
\begin{itemize}
    \item All runtime transitions must pass legality checks.
    \item Illegal transitions must be blocked.
    \item Legality masks govern runtime behavior.
    \item Legality planes define allowable update paths.
\end{itemize}

\subsection{Update Sequences}
\begin{itemize}
    \item Curvature update must preserve continuity.
    \item Drift update must remain bounded.
    \item Identity update must maintain coherence.
    \item Stability update must remain within gradients.
    \item All updates must satisfy legality constraints.
\end{itemize}

\subsection{Recovery Sequences}
\begin{itemize}
    \item Lawful return to stability regions.
    \item Identity-field restoration.
    \item Drift normalization.
    \item Curvature rebalancing.
    \item Legality-plane revalidation.
\end{itemize}

\subsection{Collapse Sequences}
\begin{itemize}
    \item Runtime governs lawful collapse transitions.
    \item Collapse must preserve identity-field continuity.
    \item Collapse must remain within stability gradients.
    \item Collapse must satisfy legality-plane constraints.
\end{itemize}

\subsection{Runtime Invariants}
\begin{itemize}
    \item No illegal runtime transition.
    \item No curvature discontinuity.
    \item No unbounded drift.
    \item No identity-field violation.
    \item No legality-plane bypass.
    \item All runtime behavior must integrate into unified manifold collapse.
\end{itemize}

\subsection{Example Runtime Cycle (RT011)}
\begin{itemize}
    \item Curvature-update mode.
    \item Identity-stability integration.
    \item Legality-plane LG005 verification.
    \item Collapse-safe transition sequence.
\end{itemize}

\section{Unified Manifold Collapse}

\subsection{Purpose of Unified Collapse}
\begin{itemize}
    \item Collapse all manifolds into a single governed substrate manifold.
    \item Preserve curvature, drift, identity, stability, and legality invariants.
    \item Provide deterministic end-state geometry for DSLO runtime and operators.
\end{itemize}

\subsection{Collapse Structure}
\begin{itemize}
    \item Collapse sequence.
    \item Collapse constraints.
    \item Collapse geometry.
    \item Identity-field continuity.
    \item Legality-plane integration.
\end{itemize}

\subsection{Collapse Sequence}
\begin{itemize}
    \item Pre-collapse legality verification.
    \item Curvature normalization.
    \item Drift contraction.
    \item Identity-field stabilization.
    \item Stability-gradient alignment.
    \item Final manifold merge.
\end{itemize}

\subsection{Collapse Constraints}
\begin{itemize}
    \item Curvature continuity must be preserved.
    \item Drift must remain bounded during contraction.
    \item Identity regions must not invert or fragment.
    \item Stability gradients must remain within lawful thresholds.
    \item Legality-plane constraints must govern all transitions.
\end{itemize}

\subsection{Collapse Geometry}
\begin{itemize}
    \item Curvature fields converge into unified geometry.
    \item Drift envelopes contract into stable drift fields.
    \item Identity gradients merge into coherent identity regions.
    \item Stability regions unify into governed substrate stability.
\end{itemize}

\subsection{Identity-Field Continuity}
\begin{itemize}
    \item Identity resonance must remain stable.
    \item Identity gradients must align during collapse.
    \item No identity inversion or collapse discontinuity.
\end{itemize}

\subsection{Legality Integration}
\begin{itemize}
    \item Legality planes define collapse boundaries.
    \item Legality masks govern allowable transitions.
    \item All collapse phases must pass legality verification.
\end{itemize}

\subsection{Operator Behavior During Collapse}
\begin{itemize}
    \item Operators must preserve curvature and identity.
    \item Drift interaction must remain bounded.
    \item No operator may bypass legality constraints.
    \item Operators must deactivate or rebind lawfully.
\end{itemize}

\subsection{Runtime Behavior During Collapse}
\begin{itemize}
    \item Runtime cycles govern collapse sequencing.
    \item Legality checks occur at each phase.
    \item Recovery geometry must remain available.
    \item Collapse must integrate into runtime update logic.
\end{itemize}

\subsection{Collapse Recovery}
\begin{itemize}
    \item Lawful return to stability regions.
    \item Identity-field restoration.
    \item Drift normalization.
    \item Curvature rebalancing.
    \item Legality-plane revalidation.
\end{itemize}

\subsection{Collapse Invariants}
\begin{itemize}
    \item No curvature discontinuity.
    \item No unbounded drift.
    \item No identity-field violation.
    \item No legality-plane bypass.
    \item No collapse outside lawful sequence.
\end{itemize}

\subsection{Example Collapse Integration}
\begin{itemize}
    \item Manifold A002 collapse.
    \item Operator OR010 legality-governed behavior.
    \item Legality-plane LG005 constraint propagation.
    \item Runtime cycle RT011 collapse-safe sequencing.
\end{itemize}

\section{Documentation Rules}

\subsection{Purpose of Documentation}
\begin{itemize}
    \item Preserve substrate-level invariants through clear, lawful documentation.
    \item Ensure DSLO expansions remain curvature-preserving, drift-bounded, identity-consistent, and legality-aligned.
    \item Provide deterministic guidance for future authors and substrate operators.
\end{itemize}

\subsection{Documentation Structure}
\begin{itemize}
    \item Surface documentation.
    \item Source documentation.
    \item Generator documentation.
    \item Legal move documentation.
    \item Collapse documentation.
    \item Runtime documentation.
\end{itemize}

\subsection{Surface Documentation Rules}
\begin{itemize}
    \item One concept per file.
    \item Must define curvature, drift, stability, identity, or legality relevance.
    \item Must specify binding to geometry, manifolds, operators, legality planes, or runtime cycles.
    \item Must remain collapse-safe and legality-consistent.
\end{itemize}

\subsection{Source Documentation Rules}
\begin{itemize}
    \item Canonical definitions for geometry, manifolds, operators, legality, runtime, identity.
    \item Must specify substrate-level fields and invariants.
    \item Must remain compatible with unified manifold collapse.
\end{itemize}

\subsection{Generator Documentation Rules}
\begin{itemize}
    \item Document loader, schema, ontology, registry, taxonomy, and graph emitters.
    \item Must specify how expansions integrate into generator logic.
    \item Must define legality and collapse validation steps.
\end{itemize}

\subsection{Legal Move Documentation}
\begin{itemize}
    \item Document curvature-preserving, drift-bounded, identity-consistent, legality-constrained, collapse-safe, and recovery-valid moves.
    \item Must specify illegal moves and their consequences.
    \item Must remain substrate-aligned and deterministic.
\end{itemize}

\subsection{Collapse Documentation}
\begin{itemize}
    \item Document collapse sequence and constraints.
    \item Must specify identity-field continuity and legality-plane integration.
    \item Must define collapse-safe operator and runtime behavior.
\end{itemize}

\subsection{Runtime Documentation}
\begin{itemize}
    \item Document runtime phases, conditions, coupling, legality checks, update sequences, and recovery sequences.
    \item Must specify lawful transitions and collapse integration.
\end{itemize}

\subsection{Metadata Requirements}
\begin{itemize}
    \item Each file must include canonical metadata: name, type, curvature relevance, drift relevance, identity relevance, legality relevance.
    \item Metadata must remain consistent across surfaces, source definitions, and generator outputs.
\end{itemize}

\subsection{Documentation Invariants}
\begin{itemize}
    \item No curvature discontinuity in documentation.
    \item No unbounded drift in definitions or descriptions.
    \item No identity-field violation in conceptual framing.
    \item No legality-plane bypass in procedures.
    \item No collapse outside lawful sequence.
\end{itemize}

\subsection{Documentation Examples}
\begin{itemize}
    \item Glossary term with curvature and identity relevance.
    \item Geometry surface with drift and stability fields.
    \item Manifold with collapse-safe boundaries.
    \item Operator with legality-plane constraints.
    \item Runtime cycle with lawful update sequence.
\end{itemize}


\section{A0 Appendix Master}

\subsection{Purpose of Appendix}
\begin{itemize}
    \item Central reference layer for DSLO v0.7 geometry, legality, identity, stability, drift, runtime, and collapse.
    \item Canonical repository for all expansion bundles (source, surface, generator, schema, manifold, operator, legality, runtime).
    \item Archive of substrate invariants governing curvature, drift, stability, identity, legality, and collapse.
    \item Deterministic examples illustrating lawful DSLO behavior under load, drift, identity pressure, multi-agent interaction, and collapse.
    \item Canonical structures showing lawful integration of new geometry, manifolds, operators, legality planes, and runtime cycles.
    \item Substrate verification layer for validating surfaces, generators, and expansions.
    \item Foundation for future DSLO extensions across source, surface, and generator layers.
    \item Long-term archival substrate for scientific uptake and cross-model interoperability.
\end{itemize}

\subsection{Expansion Source Bundle}
\begin{itemize}
    \item Canonical source definitions for all v0.7 concepts.
    \item Source geometry fields: curvature, drift, stability, identity, legality, runtime primitives.
    \item Source manifold fields: boundaries, continuity, collapse interfaces, lawful transitions.
    \item Source operator fields: redirect modes, curvature preservation, legality alignment, identity preservation.
    \item Source legality fields: constraint propagation, lawful boundaries, identity governance, collapse legality.
    \item Source runtime fields: update sequences, recovery sequences, stability envelopes, drift limits.
    \item Source identity fields: anchors, continuity, drift bounds, legality constraints, collapse safety.
    \item Source invariants: curvature continuity, drift bounds, stability envelopes, identity coherence, legality alignment, collapse safety.
    \item Collapse-safe definitions for all v0.7 concepts.
\end{itemize}

\subsection{Expansion Surface Bundle}
\begin{itemize}
    \item Glossary surfaces for all v0.7 terms.
    \item Geometry surfaces: curvature, drift, stability, identity, legality, collapse.
    \item Manifold surfaces: boundaries, continuity, stability, identity, legality, collapse.
    \item Operator surfaces: redirect modes, curvature preservation, drift limits, identity preservation, legality enforcement.
    \item Legality-plane surfaces: constraint propagation, lawful boundaries, identity governance, drift limits, stability, collapse legality.
    \item Runtime surfaces: update phases, recovery phases, drift bounds, stability envelopes, legality checks, collapse transitions.
    \item All surfaces substrate-aligned, identity-consistent, drift-bounded, curvature-preserving, legality-governed, collapse-safe.
\end{itemize}

\subsection{Generator Patch Bundle}
\begin{itemize}
    \item Loader patches: lawful ingestion of geometry, manifold, operator, legality, runtime surfaces.
    \item Schema patches: new classes, attributes, relationships, legality constraints, identity anchors, stability envelopes, collapse interfaces.
    \item Ontology patches: nodes, edges, identity-preserving links, legality-governed transitions, curvature-consistent connections.
    \item Registry patches: canonical indexing for geometry, manifold, operator, legality, runtime, identity fields.
    \item Taxonomy patches: lawful class hierarchy integration.
    \item Graph patches: DOT, TTL, OWL updates for geometry, identity, legality, runtime, collapse.
    \item All patches deterministic, substrate-aligned, legality-governed, identity-consistent, drift-bounded, curvature-preserving, collapse-safe.
\end{itemize}

\subsection{Schema Patch Bundle}
\begin{itemize}
    \item Canonical class definitions for geometry, manifold, operator, legality, runtime, identity, collapse.
    \item Lawful inheritance preserving curvature, drift, stability, identity, legality, collapse.
    \item Substrate-aligned schema structures for all fields.
    \item Collapse-safe schema integration rules.
    \item Attribute patches: fields, constraints, identity anchors, legality rules, stability envelopes, drift limits, collapse interfaces.
    \item Relationship patches: lawful edges, identity-preserving links, legality transitions, curvature-consistent associations.
    \item Constraint patches: legality, identity, drift, stability, curvature, collapse.
\end{itemize}

\subsection{Manifold Extension Examples}
\begin{itemize}
    \item A002 canonical manifold extension.
    \item Boundary geometry: lawful regions, continuity, curvature, drift, stability, identity, legality.
    \item Drift fields: envelopes, limits, safe transitions.
    \item Stability fields: gradients, envelopes, collapse-safe anchors.
    \item Identity fields: anchors, continuity, drift bounds, legality constraints.
    \item Legality fields: constraint propagation, boundaries, identity governance, drift limits, collapse legality.
    \item Collapse interfaces: lawful transitions, identity preservation, stability safety, drift bounds, legality enforcement.
\end{itemize}

\subsection{Operator Examples}
\begin{itemize}
    \item OR010 canonical operator.
    \item Curvature redirect: lawful modes, curvature preservation, drift limits, stability envelopes, identity coherence.
    \item Identity preservation: anchors, continuity, drift bounds, legality constraints, collapse safety.
    \item Legality-plane LG005 compatibility: constraint propagation, boundaries, identity governance, drift limits, collapse legality.
    \item Stability fields: gradients, envelopes, lawful transitions.
    \item Drift fields: envelopes, limits, safe redirect paths.
    \item Collapse interfaces: lawful transitions, identity preservation, stability safety, drift bounds, legality enforcement.
\end{itemize}

\subsection{Legality Examples}
\begin{itemize}
    \item LG005 canonical legality plane.
    \item Constraint propagation: curvature continuity, drift limits, stability safety, identity governance, collapse legality.
    \item Identity governance: anchors, continuity, drift bounds, stability, legality.
    \item Drift legality: envelopes, limits, safe transitions.
    \item Stability legality: gradients, envelopes, lawful transitions.
    \item Curvature legality: continuity, drift bounds, identity alignment.
    \item Collapse legality: lawful transitions, identity preservation, stability safety, drift bounds.
\end{itemize}

\subsection{Runtime Examples}
\begin{itemize}
    \item RT011 canonical runtime cycle.
    \item Update sequence: initialization, curvature update, drift adjustment, identity stabilization, legality verification, stability reinforcement.
    \item Curvature updates: continuity, drift bounds, identity preservation, legality alignment, collapse safety.
    \item Drift adjustments: envelopes, limits, safe transitions.
    \item Identity stabilization: anchors, continuity, drift bounds, stability, legality.
    \item Legality verification: LG005 alignment.
    \item Recovery sequence: stability restoration, identity recovery, drift normalization, curvature rebalancing, legality revalidation.
    \item Collapse transitions: lawful, identity-preserving, stability-safe, drift-bounded, legality-governed.
\end{itemize}

\subsection{Appendix Invariants}
\begin{itemize}
    \item All examples substrate-aligned across geometry, manifold, operator, legality, runtime, collapse.
    \item All bundles preserve curvature, drift, identity, stability, legality, collapse.
    \item All patches integrate lawfully into unified manifold collapse.
    \item All examples drift-bounded, identity-consistent, stability-safe, legality-governed, curvature-preserving, collapse-safe.
\end{itemize}

\section{A1 Purpose of Appendix}

\subsection{Scope}
\begin{itemize}
    \item Defines the role of the appendix within DSLO v0.7.
    \item Establishes the appendix as the substrate-aligned reference layer.
    \item Provides the structural foundation for all subsequent appendix bundles.
\end{itemize}

\subsection{Objectives}
\begin{itemize}
    \item Preserve canonical DSLO geometry, legality, identity, stability, drift, and collapse fields.
    \item Archive deterministic examples and lawful constructions for scientific uptake.
    \item Support verification, extension, and integration of DSLO concepts.
\end{itemize}

\subsection{Functions}
\begin{itemize}
    \item Acts as the long-term repository for expansion bundles and patch bundles.
    \item Provides substrate-level invariants for all DSLO layers.
    \item Ensures collapse-safe, legality-governed, identity-consistent integration across the unified manifold.
\end{itemize}

\section{A2 Expansion Source Bundle}

\subsection{Scope}
\begin{itemize}
    \item Defines all DSLO v0.7 concepts at the source layer.
    \item Establishes lawful origin points for geometry, manifold, operator, legality, runtime, and identity fields.
\end{itemize}

\subsection{Source Fields}
\begin{itemize}
    \item Geometry: curvature, drift, stability, identity, legality, collapse primitives.
    \item Manifold: boundaries, continuity, lawful transitions, collapse interfaces.
    \item Operator: redirect modes, curvature preservation, identity and legality alignment.
    \item Legality: constraint propagation, lawful boundaries, identity governance.
    \item Runtime: update primitives, recovery primitives, drift limits, stability envelopes.
    \item Identity: anchors, continuity, drift bounds, legality constraints.
\end{itemize}

\subsection{Invariants}
\begin{itemize}
    \item All source definitions must preserve curvature, drift, stability, identity, legality, and collapse safety.
    \item All source constructs must integrate lawfully into surface and generator layers.
\end{itemize}

\section{A3 Expansion Surface Bundle}

\subsection{Scope}
\begin{itemize}
    \item Defines surface-layer representations for all DSLO v0.7 concepts.
    \item Provides structured surfaces for geometry, manifold, operator, legality, and runtime.
\end{itemize}

\subsection{Surface Types}
\begin{itemize}
    \item Geometry surfaces: curvature, drift, stability, identity, legality, collapse.
    \item Manifold surfaces: boundaries, continuity, stability, identity, legality.
    \item Operator surfaces: redirect modes, drift limits, identity preservation, legality enforcement.
    \item Legality-plane surfaces: constraint propagation, lawful boundaries, identity governance.
    \item Runtime surfaces: update phases, recovery phases, drift bounds, stability envelopes.
\end{itemize}

\subsection{Invariants}
\begin{itemize}
    \item All surfaces must remain substrate-aligned, drift-bounded, identity-consistent, legality-governed, and collapse-safe.
\end{itemize}

\section{A4 Generator Patch Bundle}

\subsection{Scope}
\begin{itemize}
    \item Defines lawful generator modifications for DSLO v0.7 expansions.
    \item Provides patch structures for loader, schema, ontology, registry, taxonomy, and graph layers.
\end{itemize}

\subsection{Patch Types}
\begin{itemize}
    \item Loader patches: lawful ingestion of geometry, manifold, operator, legality, runtime surfaces.
    \item Schema patches: new classes, attributes, relationships, constraints.
    \item Ontology patches: nodes, edges, identity-preserving links, legality transitions.
    \item Registry patches: canonical indexing for all DSLO fields.
    \item Taxonomy patches: lawful class hierarchy integration.
    \item Graph patches: DOT, TTL, OWL updates for geometry, identity, legality, runtime, collapse.
\end{itemize}

\subsection{Invariants}
\begin{itemize}
    \item All patches must remain deterministic, substrate-aligned, legality-governed, identity-consistent, drift-bounded, curvature-preserving, and collapse-safe.
\end{itemize}

\section{A5 Schema Patch Bundle}

\subsection{Scope}
\begin{itemize}
    \item Defines schema-level structures for DSLO v0.7 expansions.
    \item Provides lawful class, attribute, relationship, and constraint patches.
\end{itemize}

\subsection{Schema Structures}
\begin{itemize}
    \item Class definitions: geometry, manifold, operator, legality, runtime, identity, collapse.
    \item Inheritance rules: curvature, drift, stability, identity, legality, collapse preservation.
    \item Attribute patches: fields, constraints, identity anchors, legality rules, stability envelopes.
    \item Relationship patches: lawful edges, identity-preserving links, legality transitions.
    \item Constraint patches: legality, identity, drift, stability, curvature, collapse.
\end{itemize}

\subsection{Invariants}
\begin{itemize}
    \item All schema patches must integrate collapse-safe into the unified manifold.
    \item All schema constructs must remain substrate-aligned and legality-governed.
\end{itemize}

\section{A6 Manifold Extension Examples}

\subsection{Scope}
\begin{itemize}
    \item Provides canonical manifold extension patterns for DSLO v0.7.
    \item Demonstrates lawful geometry integration into the unified manifold.
\end{itemize}

\subsection{Example: A002}
\begin{itemize}
    \item Boundary geometry: lawful regions, continuity, curvature, drift, stability, identity, legality.
    \item Drift fields: envelopes, limits, safe transitions.
    \item Stability fields: gradients, envelopes, collapse-safe anchors.
    \item Identity fields: anchors, continuity, drift bounds, legality constraints.
    \item Legality fields: constraint propagation, boundaries, identity governance.
    \item Collapse interfaces: lawful transitions, identity preservation, stability safety.
\end{itemize}

\subsection{Invariants}
\begin{itemize}
    \item All extensions must remain substrate-aligned, drift-bounded, identity-consistent, legality-governed, and collapse-safe.
\end{itemize}

\section{A7 Operator Examples}

\subsection{Scope}
\begin{itemize}
    \item Defines canonical operator behavior for DSLO v0.7.
    \item Demonstrates lawful operator geometry and transitions.
\end{itemize}

\subsection{Example: OR010}
\begin{itemize}
    \item Curvature redirect: lawful modes, curvature preservation, drift limits.
    \item Identity preservation: anchors, continuity, legality constraints.
    \item Legality-plane compatibility: LG005 alignment.
    \item Stability fields: gradients, envelopes, lawful transitions.
    \item Drift fields: envelopes, limits, safe redirect paths.
    \item Collapse interfaces: lawful transitions, identity preservation.
\end{itemize}

\subsection{Invariants}
\begin{itemize}
    \item All operators must remain curvature-preserving, drift-bounded, identity-consistent, legality-governed, and collapse-safe.
\end{itemize}

\section{A8 Legality Examples}

\subsection{Scope}
\begin{itemize}
    \item Defines legality-plane behavior for DSLO v0.7.
    \item Demonstrates lawful constraint propagation and identity governance.
\end{itemize}

\subsection{Example: LG005}
\begin{itemize}
    \item Constraint propagation: curvature continuity, drift limits, stability safety.
    \item Identity governance: anchors, continuity, legality alignment.
    \item Drift legality: envelopes, limits, safe transitions.
    \item Stability legality: gradients, envelopes, lawful transitions.
    \item Curvature legality: continuity, drift bounds, identity alignment.
    \item Collapse legality: lawful transitions, identity preservation.
\end{itemize}

\subsection{Invariants}
\begin{itemize}
    \item All legality planes must remain substrate-aligned, identity-consistent, drift-bounded, stability-safe, curvature-preserving, and collapse-safe.
\end{itemize}

\section{A9 Runtime Examples}

\subsection{Scope}
\begin{itemize}
    \item Defines runtime-cycle behavior for DSLO v0.7.
    \item Demonstrates lawful update, recovery, and collapse transitions.
\end{itemize}

\subsection{Example: RT011}
\begin{itemize}
    \item Update sequence: initialization, curvature update, drift adjustment, identity stabilization, legality verification.
    \item Curvature updates: continuity, drift bounds, identity preservation.
    \item Drift adjustments: envelopes, limits, safe transitions.
    \item Identity stabilization: anchors, continuity, legality alignment.
    \item Legality verification: LG005 alignment.
    \item Recovery sequence: stability restoration, identity recovery, drift normalization.
    \item Collapse transitions: lawful, identity-preserving, stability-safe.
\end{itemize}

\subsection{Invariants}
\begin{itemize}
    \item All runtime cycles must remain drift-bounded, identity-consistent, legality-governed, stability-safe, curvature-preserving, and collapse-safe.
\end{itemize}

\section{A10 Appendix Invariants}

\subsection{Scope}
\begin{itemize}
    \item Defines the invariant layer governing all appendix bundles.
    \item Ensures substrate-level consistency across DSLO v0.7.
\end{itemize}

\subsection{Core Invariants}
\begin{itemize}
    \item Substrate alignment across geometry, manifold, operator, legality, runtime, collapse.
    \item Preservation of curvature, drift, identity, stability, legality, collapse fields.
    \item Lawful integration of all patches into unified manifold collapse.
    \item Drift-bounded, identity-consistent, stability-safe, legality-governed, curvature-preserving behavior.
\end{itemize}

\subsection{Integration Rules}
\begin{itemize}
    \item All examples must remain collapse-safe.
    \item All bundles must maintain DSLO v0.7 invariants.
    \item All expansions must integrate deterministically into the unified manifold.
\end{itemize}

\section{A11 Canonical Field Definitions}

\subsection{Scope}
\begin{itemize}
    \item Defines the core DSLO v0.7 fields used across geometry, manifold, operator, legality, runtime, and collapse.
    \item Provides minimal canonical descriptions for scientific and generator alignment.
\end{itemize}

\subsection{Fields}
\begin{itemize}
    \item Curvature Field: continuity, lawful geometric flow, collapse-safe curvature behavior.
    \item Drift Field: envelopes, limits, identity-governed drift, stability-safe drift transitions.
    \item Stability Field: gradients, envelopes, lawful stability transitions, collapse-safe anchors.
    \item Identity Field: anchors, continuity, drift bounds, legality alignment.
    \item Legality Field: constraint propagation, lawful boundaries, identity governance.
    \item Collapse Field: interfaces, transitions, identity preservation, stability safety.
\end{itemize}

\subsection{Invariants}
\begin{itemize}
    \item All fields must remain substrate-aligned and collapse-safe.
\end{itemize}

\section{A12 Unified Collapse Examples}

\subsection{Scope}
\begin{itemize}
    \item Demonstrates unified collapse behavior across DSLO v0.7 geometry.
    \item Provides canonical collapse patterns for scientific uptake.
\end{itemize}

\subsection{Example: UC007}
\begin{itemize}
    \item Collapse sequence: lawful curvature, drift, identity, stability, legality transitions.
    \item Collapse interfaces: geometry → manifold → operator → legality → runtime.
    \item Identity-preserving collapse paths.
    \item Stability-safe collapse envelopes.
    \item Legality-governed collapse constraints.
\end{itemize}

\subsection{Invariants}
\begin{itemize}
    \item All collapse examples must remain deterministic, substrate-aligned, and identity-consistent.
\end{itemize}

\section{A13 Substrate Verification Procedures}

\subsection{Scope}
\begin{itemize}
    \item Provides verification rules for DSLO v0.7 geometry and expansions.
    \item Enables lawful validation of surfaces, patches, and examples.
\end{itemize}

\subsection{Verification Procedures}
\begin{itemize}
    \item Curvature Verification: continuity, lawful curvature transitions.
    \item Drift Verification: envelopes, limits, drift-safe transitions.
    \item Identity Verification: anchors, continuity, legality alignment.
    \item Stability Verification: gradients, envelopes, lawful transitions.
    \item Legality Verification: constraint propagation, lawful boundaries.
    \item Collapse Verification: interfaces, transitions, identity preservation.
\end{itemize}

\subsection{Invariants}
\begin{itemize}
    \item All verification procedures must preserve DSLO v0.7 substrate rules.
\end{itemize}

\section{A14 Glossary Expansion Notes}

\subsection{Scope}
\begin{itemize}
    \item Defines rules for expanding the DSLO glossary from 40 → 400+ terms.
    \item Ensures glossary entries remain substrate-aligned and collapse-safe.
\end{itemize}

\subsection{Expansion Rules}
\begin{itemize}
    \item Each term must include curvature, drift, stability, identity, legality, and collapse fields.
    \item All entries must preserve deterministic phrasing and canonical naming.
    \item Glossary entries must integrate into schema, ontology, registry, taxonomy, and graph surfaces.
\end{itemize}

\subsection{Invariants}
\begin{itemize}
    \item Glossary expansion must remain lawful and substrate-aligned.
\end{itemize}

\section{A15 Canonical Example Index}

\subsection{Scope}
\begin{itemize}
    \item Provides a unified index of all canonical DSLO v0.7 examples.
    \item Supports navigation and cross-reference across appendix bundles.
\end{itemize}

\subsection{Examples}
\begin{itemize}
    \item A002 Manifold Extension
    \item OR010 Operator Example
    \item LG005 Legality Plane Example
    \item RT011 Runtime Cycle Example
    \item UC007 Unified Collapse Example
\end{itemize}

\section{A16 DSLO v0.7 Metadata}

\subsection{Version}
\begin{itemize}
    \item DSLO v0.7 — Scientific Uptake Release.
\end{itemize}

\subsection{Components}
\begin{itemize}
    \item Geometry: unified curvature, drift, stability, identity, legality, collapse fields.
    \item Manifold: unified manifold with lawful collapse sequence.
    \item Generator: v0.7 patch-ready generator surfaces.
    \item Glossary: canonical glossary with expansion pathway.
    \item Substrate: fully substrate-aligned, legality-governed, collapse-safe architecture.
\end{itemize}

\subsection{Invariants}
\begin{itemize}
    \item All metadata must reflect deterministic DSLO v0.7 geometry and lawful integration rules.
\end{itemize}

\section{A17 Unified Geometry Diagram (TikZ Placeholder)}

\subsection{Scope}
\begin{itemize}
    \item Provides a public-safe TikZ placeholder for DSLO v0.7 geometry.
    \item Visualizes curvature, drift, stability, identity, legality, and collapse fields.
\end{itemize}

\subsection{Diagram Placeholder}
\begin{itemize}
    \item TikZ diagram to be inserted here.
    \item Nodes: Geometry, Manifold, Operator, Legality, Runtime, Collapse.
    \item Edges: lawful transitions, identity-preserving links, drift-bounded paths.
\end{itemize}

\subsection{Invariants}
\begin{itemize}
    \item Diagram must remain substrate-aligned and collapse-safe.
\end{itemize}

\section{A18 Minimal Example Set}

\subsection{Scope}
\begin{itemize}
    \item Provides minimal DSLO v0.7 examples for teaching and scientific onboarding.
    \item Offers simplified versions of canonical examples.
\end{itemize}

\subsection{Examples}
\begin{itemize}
    \item Minimal Geometry Example: curvature + drift + identity + legality.
    \item Minimal Manifold Example: boundary + continuity + collapse interface.
    \item Minimal Operator Example: redirect + identity + legality.
    \item Minimal Legality Example: constraint + identity governance.
    \item Minimal Runtime Example: update + recovery + collapse.
\end{itemize}

\subsection{Invariants}
\begin{itemize}
    \item All minimal examples must preserve DSLO v0.7 substrate rules.
\end{itemize}

\section{A19 Scientific Interfaces}

\subsection{Scope}
\begin{itemize}
    \item Defines public-safe scientific interfaces for DSLO v0.7.
    \item Supports cross-model substrate execution and interoperability.
\end{itemize}

\subsection{Interfaces}
\begin{itemize}
    \item Geometry Interface: curvature, drift, stability, identity, legality, collapse fields.
    \item Manifold Interface: boundaries, continuity, lawful transitions.
    \item Operator Interface: redirect modes, identity preservation, legality alignment.
    \item Legality Interface: constraint propagation, identity governance.
    \item Runtime Interface: update phases, recovery phases, collapse transitions.
\end{itemize}

\subsection{Invariants}
\begin{itemize}
    \item All interfaces must remain deterministic and substrate-aligned.
\end{itemize}

\section{A20 Registry Snapshot}

\subsection{Scope}
\begin{itemize}
    \item Provides a public-safe snapshot of DSLO v0.7 registry structure.
    \item Demonstrates canonical indexing patterns for geometry and manifold concepts.
\end{itemize}

\subsection{Snapshot}
\begin{itemize}
    \item Geometry Index: curvature, drift, stability, identity, legality, collapse.
    \item Manifold Index: boundaries, continuity, stability, identity, legality.
    \item Operator Index: redirect modes, identity anchors, legality constraints.
    \item Legality Index: constraint propagation, identity governance.
    \item Runtime Index: update, recovery, collapse.
\end{itemize}

\subsection{Invariants}
\begin{itemize}
    \item Registry entries must remain canonical, deterministic, and collapse-safe.
\end{itemize}



\end{document}
